\chapter{Chronic Cough}

\section*{Definition}
A chronic cough is a cough lasting more than eight weeks in adults or more than four weeks in children. It may reflect airway disease, reflux, upper-airway inflammation, medication effects, infection, or lung disease.

\section*{When to See a Doctor}
Seek urgent evaluation for Coughing blood, breathlessness, chest pain, fever, weight loss, night sweats, fainting, or low oxygen. Arrange routine or prompt assessment for persistent, recurrent, unexplained, or function-limiting symptoms.

\section*{How It Develops}
The symptom reflects irritation, inflammation, obstruction, altered secretion, infection, or dysfunction in the relevant organ system. The pattern, duration, associated symptoms, exposures, and examination findings determine whether it is self-limited or requires investigation.

\section*{Causes}
\begin{itemize}
  \item Asthma, upper-airway cough syndrome, gastroesophageal reflux, smoking, ACE inhibitors, chronic infection, bronchiectasis, interstitial disease, and malignancy.
  \item Medication effects, environmental exposures, and chronic disease may contribute.
  \item Less common but important causes include malignancy, vascular disease, or organ failure when symptoms persist or progress.
\end{itemize}

\section*{Related Symptoms}
\begin{itemize}
  \item Pain, fever, fatigue, nausea, or changes in normal organ function
  \item Bleeding, weight change, shortness of breath, weakness, or neurologic symptoms when a serious cause is present
\end{itemize}

\section*{Medical Evaluation}
\begin{itemize}
  \item History addresses onset, duration, severity, triggers, exposures, medications, medical conditions, pregnancy status when relevant, and warning symptoms.
  \item Examination focuses on vital signs and the organ systems suggested by the symptom.
  \item Testing may include blood or urine studies, cultures, imaging, physiologic testing, or specialist examination according to the clinical pattern.
\end{itemize}

\section*{Treatment Options}
Treatment is directed at the cause and may include supportive care, targeted medication, treatment of infection or inflammation, correction of metabolic disease, or specialist procedures. Persistent symptoms require reassessment rather than indefinite symptomatic treatment.

\section*{Self-Care and Home Management}
Avoid known triggers, maintain reasonable hydration, record timing and associated symptoms, and use only treatments appropriate for the suspected cause. Do not use leftover antibiotics or delay urgent care because symptoms temporarily improve.

\section*{References}
\begin{thebibliography}{99}
\bibitem{chronic_cough:accp} Irwin RS, Baumann MH, Bolser DC, et al. Diagnosis and management of cough executive summary: ACCP evidence-based clinical practice guidelines. \textit{Chest}. 2006;129(suppl 1):1S--23S.
\bibitem{chronic_cough:ers} Morice AH, Millqvist E, Bieksiene K, et al. ERS guidelines on the diagnosis and treatment of chronic cough in adults and children. \textit{Eur Respir J}. 2020;55:1901136.
\bibitem{chronic_cough:harrison} Jameson JL, et al. \textit{Harrison's Principles of Internal Medicine}. 21st ed. McGraw-Hill; 2022.
\end{thebibliography}
